\documentclass{article}
\usepackage{graphicx} % Required for inserting images
\usepackage{amsmath}
\title{Lesson 3 Homework - Proof Writing}
\author{William Wang}
\date{September 2025}

\begin{document}

\maketitle

\section{}
We are give that \(8^x = 27\) and we want to find what \(4^{2x-3}\) is. We can first change them both into the same base.
\[8^x = 27 \implies 2^{3x} = 3^3 \text{ and } 4^{x-3} = 2^{4x-6}\]
We can get \( 2^{4x-6}\) by starting with \(2^{3x}\), multiplying by \(2^x\) and dividing by \(2^6\). We can get \(2^x\) by taking the cube root of each side in the first equation, giving us
\[2^x = 3\] Hence, 
\[2^{4x-6 }= 2^{3x} \cdot 2^{x}  \div 2^6  = 27 \cdot 3 \div 64 = \frac{81}{64}\]
\section{}
We start with \(\log_{a^n}b\), by the change of base formula, we have 
\[\log_{a^n} b = \frac{1}{\log_b a^n}\]
Then by the power rule for exponents, we have
\[\frac{1}{\log_b a^n} = \frac{1}{n\log_b a} = \frac{1}{n}\cdot \frac{1}{\log_b a} = \frac{1}{n} \log_a b\]
Hence, 
\[\log_{a^n} b = \frac{1}{n} \log_a b\]
\section{}
To find the point that they meet, we can solve the simultaneous equation
\[y = \log_2 (x+3)\]
\[y = 3-\log_2 (x-1)\]

\[\implies \log_2(x+3) = 3-\log_2(x-1)\]
\[\implies \log_2(x+3) = \log_2 8 - \log_2(x-1)\]
\[\implies \log_2(x+3)=\log_2(\frac{8}{x-1})\]
Since the logarithmic function has monotonicity, we have
\[x+3 = \frac{8}{x-1} \implies x^2+2x-3 = 8 \implies x^2+2x-11 = 0\]
Notice the solution for the quadratic has to be positive since from the second equation we know that \(x > 1\). 
\[\implies x = \frac{-2+\sqrt{48}}{2} = -1 + 2\sqrt{3}\]
Then we have $y = \log_2(2+2\sqrt{3} = 1 + \log_2 (1+\sqrt{3})$. Hence the point is 
\[(-1+2\sqrt{3}, 1+\log_2(1+\sqrt{3}))\]


\end{document}
